\section{Marco Teórico}
\label{ch:marco_teorico}

Para el desarrollo del Sistema de Apoyo a la Generación de Acuerdos Judiciales (SAGAI), es indispensable establecer los fundamentos teóricos y tecnológicos que soportan la arquitectura del software. A continuación, se describen los conceptos clave, su funcionamiento técnico y su pertinencia en el entorno legal.

\subsection{Transformación Digital y LegalTech}

\subsubsection{Conceptos fundamentales y evolución}
El término \textit{LegalTech} se refiere a la intersección entre la tecnología de la información y la prestación de servicios legales. En el sector público, esto implica la digitalización de expedientes, la automatización de procesos administrativos y la implementación de sistemas de gestión documental inteligente. La modernización tecnológica representa un cambio paradigmático en la administración de justicia, impulsando la adopción de herramientas predictivas y de soporte institucional que facilitan el acceso equitativo a la justicia \cite{Frontiers2025}. 

\subsubsection{Impacto en la eficiencia de los tribunales}
La integración de tecnologías de la Industria 4.0 en la administración de los tribunales ha demostrado ser fundamental para resolver crisis operativas estructurales \cite{Bhatt2024}. La acumulación de casos (\textit{backlog}) y las ineficiencias en los flujos de trabajo generan retrasos sistémicos. Implementar sistemas de gestión y análisis de datos optimiza las operaciones de servicio, permitiendo una asignación eficiente del tiempo y recursos procesales \cite{Bakshi2024}. Además, los Sistemas de Gestión de Casos Judiciales (CCMS) han comprobado mejorar drásticamente la administración documental y el intercambio de información, sentando la base para ecosistemas digitales eficientes \cite{CourtCaseManagement2026}.

\subsection{Procesamiento Digital de Documentos Jurídicos}

\subsubsection{Gestión de expedientes electrónicos}
La digitalización de los entornos judiciales exige métodos robustos para el almacenamiento y gestión de la información. Más allá de la simple conversión de papel a formato digital, el procesamiento de documentos jurídicos modernos requiere infraestructuras capaces de organizar semánticamente el inmenso acervo judicial.

\subsubsection{Enfoque centrado en entidades}
En este contexto, investigaciones recientes proponen un enfoque centrado en entidades para la gestión de decisiones judiciales. Mediante el pre-procesamiento iterativo con servicios de Procesamiento de Lenguaje Natural, es posible extraer automáticamente entidades legales clave (tales como partes involucradas, jueces y estatutos citados). Este tipo de anotación semántica facilita enormemente la gestión del repositorio digital, dotando de estructura y capacidad de búsqueda a documentos que de otra manera serían archivos estáticos \cite{Bellandi2024}.

\subsection{Reconocimiento Óptico de Caracteres (OCR)}

\subsubsection{Principios de extracción de texto}
El Reconocimiento Óptico de Caracteres (OCR) es una tecnología de visión computacional orientada a la detección y conversión de patrones de pixeles en caracteres alfanuméricos codificados por máquina. Permite transformar diferentes tipos de documentos estáticos, como imágenes escaneadas o archivos PDF fotográficos, en datos de texto editable.

\subsubsection{Arquitectura y Tesseract OCR}
Los avances recientes en OCR integran modelos de aprendizaje profundo que optimizan el análisis de disposición (\textit{Document Layout Analysis}) y la detección precisa de líneas de texto \cite{Fateh2024}. El uso de motores de OCR de código abierto, como Tesseract, ha sido ampliamente evaluado en el entorno jurídico. Estudios enfocados en documentos legales digitalizados demuestran que la integración del OCR con modelos modernos de visión mejora sustancialmente la extracción de contenido. Esto valida a Tesseract OCR como una herramienta altamente viable y efectiva para la implementación de \textit{pipelines} de digitalización judicial de bajo costo \cite{OCRLegal2025}.

% --- INSERTA AQUÍ TU IMAGEN DE OCR ---
\begin{figure}[H]
    \centering
    % Descomenta la siguiente línea y pon el nombre exacto de la imagen que subas a tu carpeta img
    % \includegraphics[width=0.7\textwidth]{img/tu_imagen_ocr.png} 
    \rule{0.7\textwidth}{4cm} % Cuadro gris temporal (bórralo cuando pongas tu imagen)
    \caption{Ejemplo del proceso de análisis de disposición y extracción de texto mediante OCR.}
    \label{fig:ejemplo_ocr}
\end{figure}

\subsection{Procesamiento de Lenguaje Natural (NLP) en Textos Legales}

\subsubsection{Fundamentos del NLP en el dominio jurídico}
El Procesamiento de Lenguaje Natural (NLP) es una rama de la inteligencia artificial que dota a los sistemas informáticos de la capacidad de comprender, interpretar y procesar el lenguaje humano. La aplicación del NLP en contextos legales presenta desafíos únicos debido a características estructurales específicas: documentos de extensión considerable, léxico altamente especializado y formatos no estructurados \cite{ACMComputing2024}.

\subsubsection{Tareas, modelos y desafíos}
Las investigaciones enfocadas en el dominio legal demuestran que el NLP permite automatizar tareas complejas como la clasificación de documentos procesales y la búsqueda semántica de información \cite{Ariai2025}. Estas herramientas metodológicas proporcionan la base algorítmica para que el sistema informático estructure la información contenida en las promociones.

\subsection{Modelos de Lenguaje Grande (LLMs) y Arquitectura Transformer}

\subsubsection{Arquitectura Transformer y autoatención}
En años recientes, el NLP ha sido revolucionado por los Modelos de Lenguaje Grande (LLMs). Estas arquitecturas de redes neuronales se fundamentan en el modelo \textit{Transformer} y su mecanismo de autoatención (\textit{self-attention}). Este mecanismo permite que cada posición en una secuencia de entrada establezca relaciones ponderadas con todas las demás palabras, capturando dependencias a largo plazo y el contexto global del documento \cite{Chen2025}.

\subsubsection{Progreso y aplicación en entornos legales}
La comunidad científica ha validado la viabilidad técnica de los LLMs para ejecutar tareas de análisis documental y síntesis de expedientes, superando las limitaciones tradicionales de procesamiento mediante técnicas de \textit{prompting} \cite{SurveyLLMLegal2026}.

% --- INSERTA AQUÍ TU IMAGEN DE LLM O IA ---
\begin{figure}[H]
    \centering
    % \includegraphics[width=0.6\textwidth]{img/tu_imagen_llm.png} 
    \rule{0.6\textwidth}{4cm} % Cuadro gris temporal (bórralo cuando pongas tu imagen)
    \caption{Representación del mecanismo de autoatención en la arquitectura Transformer.}
    \label{fig:arquitectura_llm}
\end{figure}

\subsection{Inteligencia Artificial Generativa para Documentos Legales}

\subsubsection{Generación de documentos y estructuración normativa}
Los modelos de inteligencia artificial generativa han demostrado capacidades transformativas en la redacción de documentos formales. En sistemas jurídicos inteligentes, técnicas como el uso de Grafos de Conocimiento permiten estructurar la normatividad procesal y utilizar relaciones semánticas para generar borradores más completos y coherentes que los sistemas basados en reglas rígidas \cite{KnowledgeGraph2024}. 

\subsubsection{Redacción asistida y mitigación de riesgos legales}
El uso de IA generativa eleva significativamente la productividad en la elaboración de acuerdos preliminares y la resolución de conflictos \cite{ChienKim2026}. Adicionalmente, el desarrollo de \textit{benchmarks} especializados como ContractEval demuestra la capacidad de los LLMs para identificar riesgos legales a nivel de cláusula, garantizando que el texto generado cumpla con los estándares de precisión requeridos \cite{ContractEval2025}.

\subsection{Generación Aumentada por Recuperación (RAG)}

\subsubsection{Contextualización de Modelos Generativos}
La Generación Aumentada por Recuperación (RAG, por sus siglas en inglés) es una arquitectura que integra modelos de lenguaje con bases de conocimiento externo y normatividad jurídica. Esta técnica es fundamental para anclar las respuestas del modelo a textos legales reales, mejorando la precisión y evitando la generación de información ficticia, un problema comúnmente conocido como alucinaciones.

\subsubsection{Implementación en la redacción asistida}
Plataformas aplicadas en el estado del arte, como el sistema JusBuild, demuestran la eficacia de RAG en la generación de documentos legales. Al combinar la recuperación semántica de precedentes similares con las sugerencias generadas por LLMs, el sistema asiste en la redacción de acuerdos. Esto permite que el personal judicial mantenga el control sobre el documento, reduciendo drásticamente las alucinaciones en contextos de jurisprudencia real \cite{JusBuild2025}.

\subsection{Inteligencia Artificial Explicable (XAI) y Ética Judicial}

\subsubsection{Transparencia y explicabilidad algorítmica}
La adopción de sistemas de IA en la impartición de justicia exige evitar el efecto de "caja negra". La Inteligencia Artificial Explicable (XAI) identifica la necesidad de proporcionar respuestas transparentes y sensibles al contexto legal, previniendo sesgos algorítmicos y asegurando que las sugerencias del sistema estén debidamente fundamentadas en la ley \cite{ICIT2024}.

\subsubsection{Percepción pública y el enfoque Human-in-the-Loop}
Estudios sobre la percepción ciudadana destacan que el uso de IA en tribunales es aceptado únicamente si se mantiene la supervisión humana, protegiendo la equidad procesal \cite{PublicPerceptionsAI2025}. La implementación exitosa de casos piloto, como el uso de asistentes virtuales por jueces en Colombia, valida la utilidad de la IA, pero subraya la responsabilidad indelegable del servidor público \cite{UnveilingAI2024}. Por ello, SAGAI se diseña como una herramienta de apoyo (\textit{Human-in-the-Loop}), cuyo propósito es aliviar el estrés judicial crónico (\textit{burnout}) causado por el trabajo manual intensivo \cite{Fine2024}, sin sustituir la toma de decisiones.

\subsection{Arquitectura Web Cliente-Servidor}

\subsubsection{Principios del modelo y flujo de interacción}
La arquitectura Cliente-Servidor constituye el paradigma fundamental para el desarrollo de la plataforma web. Esta arquitectura facilita la separación de responsabilidades, desacoplando la interfaz de usuario que se ejecuta en los navegadores (\textit{Frontend}) de la lógica de procesamiento masivo y bases de datos centralizadas en los servidores (\textit{Backend}) \cite{Nyabuto2024}. 

\subsubsection{Justificación de la arquitectura y uso de FastAPI}
Para SAGAI, este modelo permite un rendimiento optimizado al distribuir la carga computacional. El Cliente interactúa recibiendo el documento digitalizado, mientras el Servidor, apoyado en marcos de trabajo de alto rendimiento para el desarrollo de APIs RESTful como FastAPI, orquesta el proceso intensivo del OCR local y gestiona las peticiones seguras hacia el modelo de lenguaje, retornando al usuario una interfaz rápida, escalable e intuitiva.